\chapter{Related Work and Background}

\section*{Shared growth models}
Top height tracks upper-canopy development and serves as a standard metric of forest growth and site productivity. Multi-temporal airborne LiDAR surveys now allow top height to be estimated consistently across broad landscapes, though uncertainty persists in young stands, open canopies, and edge boundaries \cite{ForestResearch2026Sitka_C1_SS, suarezminguez2008Practical_C1_SS}

The Chapman--Richards model describes height trajectories using a flexible sigmoid curve parametrised by an asymptotic ceiling ($y_{\max}$), an intrinsic growth rate ($k$), and a shape exponent ($p$). While a single population curve offers a useful baseline, it cannot capture varying local site conditions and management histories \cite{richards1959Flexible_C2_CR, pommerening2015Methods_C2_CR, mason2011Sitka_C1_SS}

Site-conditioned models address this limitation by allowing environmental predictors to adjust curve parameters. However, fitting multiple parameters from short observation series can cause parameter trade-off's, as different combinations of $y_{\max}$, $k$, and $p$ yield similar height predictions over narrow age windows \cite{scolforo2013Dominant_C2_CR, socha2021Regional_C2_CR, rennolls1995Forest_C2_CR}. Mean residual departure from a shared curve and departure in fitted asymptotes capture related yet distinct aspects of productivity: the former tracks whether a stand consistently sits above or below the population mean, whereas the latter reflects site-specific height ceilings while holding rate and shape fixed.

\begin{figure}[!htbp]
    \centering
    \includegraphics[width=0.65\textwidth]
    {figures/fig_background/CR-curves-yeildclass.png}
    \caption{Chapman-Richards growth curves coloured by General Yield Class (GYC).}
    \label{fig:chapman-richards-illustrative}
\end{figure}

\paragraph*{Environmental controls}
Sitka spruce productivity responds strongly to elevation, wind exposure, soil moisture, and soil fertility. In upland northern Britain, General Yield Class falls by roughly $3.2\text{--}4.0\,\mathrm{m^3\,ha^{-1}\,yr^{-1}}$ per 100\,m rise in elevation, driven by lower temperatures and increased wind exposure \cite{worrell1990Productivity1_C1_SS}. Studies of Irish plantations similarly show peak productivity on moist, nutrient-rich sites, with sharp declines on drought-prone or waterlogged soils \cite{farrelly2011Sitka_C1_SS}.

Landscape-scale modelling integrates forest inventory data with spatial layers. Early UK research combined sub-compartment records with terrain and soil maps using principal component analysis and linear regression to predict yield classes \cite{bateman1998Using_C1_SS}. More recent studies pair LiDAR metrics with climate, topographic, and edaphic rasters. Because predictor importance varies across spatial resolutions and extents \cite{fricker2019More_C3_LID}, and coarse soil layers often correlate poorly with local soil conditions \cite{miller2024Assessing_C2_CR}, the literature confirms broad environmental controls on growth without establishing a universal hierarchy of drivers across scales.

\paragraph*{Global attribution}
Global models assume a single environment--growth relationship across the entire study area. Because findings depend heavily on assumptions regarding linearity, interactions, and collinearity, comparing diverse model families helps determine whether an observed ecological association is robust or an artifact of model structure.

Linear mixed-effects models (LMMs) isolate population-level fixed effects from grouped, hierarchically structured site variation \cite{harrison2018Brief_LMM, socha2016Assessment_C2_CR}. However, their fixed effects remain globally linear and additive unless interaction terms are explicitly parameterised. Elastic Net applies combined $L_1$ variable selection and $L_2$ shrinkage to handle collinear predictors, though its estimates remain constrained by linear functional assumptions \cite{zou2005Regularization_C4_NN}.

Gradient-boosted decision trees capture complex, non-linear environmental relationships and high-order feature interactions \cite{puliti2023Tree_C2_CR, chen2016XGBoost_C4_NN}. When paired with SHAP, tree models reveal non-linear thresholds that linear models miss \cite{lundberg2017Unified_C7_XAI}. But SHAP values reflect feature attribution within the model rather than biological causality, and importance can be partitioned arbitrarily among correlated predictors \cite{aas2021Explaining_C7_XAI}. Few studies systematically benchmark linear, mixed-effects, and tree-based attribution models on the exact same growth-departure targets.

\paragraph*{Spatially varying attribution}
Spatial autocorrelation describes spatial clustering in observations or model residuals, whereas spatial non-stationarity refers to geographic shifts in the underlying process itself. Residual spatial autocorrelation can signal omitted environmental predictors, incorrect model functional form, boundary edge effects, or genuine non-stationarity. 

Geographically Weighted Regression (GWR) models spatial variation by fitting local linear regressions weighted by proximity, producing interpretable local coefficient maps \cite{brunsdon1996Geographically_C6_GWR, wheeler2005Multicollinearity_C6_GWR}. Multiscale GWR (MGWR) further relaxes this by allowing different covariates to operate at distinct spatial bandwidths [24].

Geographically Neural Network-Weighted Regression (GNNWR) replaces predefined distance kernels with a learned neural weighting sub-network \cite{du2020Geographically_C6_GWR}. While this flexible weighting captures non-linear spatial dependencies, it reduces transparency and requires strict spatial validation. 

\paragraph*{Physics-guided prediction}
% TODO: karpatne2017Theory_C4_NN confirmed in citations.md; raissi2019Physics_C4_NN,
% karniadakis2021Physics_C4_NN, willard2022Integrating_C4_NN did not match a grep check --
% possibly a formatting mismatch (wang2021Understanding_C4_NN also didn't match, despite
% being already live-cited two paragraphs below), possibly a real gap. Confirm before final
% submission; not re-checked further per instruction to stop searching for now.
\textbf{PINNs} and hybrid models incorporate scientific relationships into the training objective \cite{karpatne2017Theory_C4_NN, raissi2019Physics_C4_NN, karniadakis2021Physics_C4_NN, willard2022Integrating_C4_NN}. A common loss formulation is:
\begin{equation}
\mathcal{L}_{\text{total}} = \mathcal{L}_{\text{data}} + \lambda\mathcal{L}_{\text{process}},
\label{eq:pinn-loss-general}
\end{equation}
where $\mathcal{L}_{\text{data}}$ measures prediction error, $\mathcal{L}_{\text{process}}$ measures disagreement with the selected process relationship, and $\lambda$ balances the two objectives.

% TODO: wesselkamp2024Process_C4_NN not confirmed in citations.md -- same caveat as above.
Process knowledge can enter neural models in several ways \cite{wesselkamp2024Process_C4_NN}:
\begin{table}[ht]
\centering
\footnotesize
\begin{tabular}{ll}
\hline
\textbf{Integration Strategy} & \textbf{Structural Mechanism} \\
\hline
\emph{Bias correction} & Network learns systematic residuals of process model outputs. \\
\emph{Parallel physics} & Process and neural predictions are combined additively. \\
\emph{Physics regularisation} & Unconstrained neural predictor is penalised for equation disagreement. \\
\emph{Domain adaptation} & Pre-training on process simulations before empirical fine-tuning. \\
\emph{Physics embedding} & Differentiable process models embedded with parameter sub-networks. \\
\hline
\end{tabular}
\label{tab:pinn-strategies}
\end{table}

In forestry applications, process-guided networks have integrated biological growth, competition, and mortality models \cite{pichler2025Inferring_C4_NN}. Conditional and variable-coefficient PINN frameworks combine physics regularisation with physics embedding, letting environmental covariates drive growth-curve parameters rather than treating them as fixed constants \cite{kovacs2022Conditional_C4_NN, miao2023VCPINN_C4_NN, scolforo2013Dominant_C2_CR}. Frameworks such as Forest-Informed Neural Networks integrate neural networks into more complex differential equation models covering growth, mortality, regeneration, and competition; a single Chapman--Richards curve is a simpler target, for which the ecological equation can instead be anchored with a soft penalty term \cite{pichler2025Inferring_C4_NN}.

This project's own PINN sits on one side of that combination, not both: it uses physics regularisation (a soft penalty on disagreement with the Chapman--Richards derivative) rather than physics embedding (a hard-constrained differentiable process model), and environmental covariates condition the parameter sub-networks' outputs ($y_{\max}^{(i)}$, $k^{(i)}$) rather than being concatenated directly into the prediction trunk's raw input. The combined regularisation-and-embedding approach the conditional/variable-coefficient literature above describes was not attempted here.

Adding process penalties creates a trade-off between fitting the observations and adhering to the Chapman--Richards guidance. Composite loss functions can produce imbalanced gradients, causing one loss component to dominate optimisation, while ill-conditioned loss landscapes can slow or prevent convergence \cite{wang2021Understanding_C4_NN}. % TODO: rathore2024Challenges_C4_NN is currently only a commented note-to-self in citations.md (line ~484: "Rathore et al. (2024) -- Challenges in training PINNs: A loss landscape perspective"), not a real bib entry yet -- add the entry before citing it, do not cite the key as-is.
There is also a risk that a strong penalty based on a misspecified population-level Chapman--Richards curve pulls predictions towards an unsuitable average rather than allowing the network to represent local growth differences.

\paragraph*{Residual spatial evaluation}
Standard predictive metrics (e.g., $R^2$, RMSE) evaluate average accuracy but fail to identify whether prediction errors remain geographically clustered. Residual spatial autocorrelation provides a critical diagnostic: persistent residual clustering indicates unmodeled spatial structure. Global Moran's $I$ evaluates residual autocorrelation against a spatial weights matrix, where reductions in absolute $I$ indicate weaker spatial clustering \cite{moran1950Notes_c5_eval, kim2021Predicting_C5_EVAL}. However, metric interpretation depends directly on the chosen spatial weights and distance thresholds. Moreover, global Moran's $I$ cannot detect cancelling local spatial clusters, and lower autocorrelation does not necessarily indicate a smaller overall error variance.

\section{Research Gap and Contributions}
Although environmental factors clearly drive forest growth, it remains uncertain whether identified drivers stay consistent across linear, mixed-effects, and non-linear models on the same data. Additionally, neural spatial models are rarely evaluated across held-out forest blocks, and the predictive value of physics-guided constraints remains unproven.

This work addresses these gaps through three contributions:
\begin{enumerate}
\setlength\itemsep{0pt}
\item \textbf{Global attribution (Q1):} Tests whether environmental drivers of curve departures remain consistent across Elastic Net, mixed-effects, and XGBoost models.
\item \textbf{Spatial attribution (Q2):} Tests spatially varying relationships using GNNWR and residual spatial autocorrelation under spatial cross-validation.
\item \textbf{Prediction and physics guidance (Q3):} Compares standard and neural models -- including a corrected physics-informed network -- on top-height prediction, testing which architecture performs best and why, and whether environmental signals or Chapman--Richards constraints improve that picture.
\end{enumerate}
This framework separates what drives growth deviations, whether those effects vary across space, and whether they improve prediction.
